% WAFC -- manuscrito, versao k = 1 (E5a). Template da Statistica Sinica
% (manuscript/ss-template/SS-template-bib.tex), decisao D5.
% Secoes 1 a 4; as Secoes 5 a 7 entram em E5b. As provas estao em supp_1.tex.
% Notacao: docs/notacao.md, congelada em E1.1. Os enunciados vem de
% derivations/01 a 05; o comentario acima de cada um diz de onde.
\documentclass[12pt]{article} %***
\usepackage[sectionbib]{natbib}
\usepackage{array,epsfig,fancyheadings,rotating}
\usepackage[]{hyperref}
%%%%%%%%%%%%%%%%%%%%%%%%%%%%%%%%%%%%
\usepackage{sectsty, secdot}
\sectionfont{\fontsize{12}{14pt plus.8pt minus .6pt}\selectfont}
\renewcommand{\theequation}{\thesection\arabic{equation}}
\subsectionfont{\fontsize{12}{14pt plus.8pt minus .6pt}\selectfont}
%%%%%%%%%%%%%%%%%%%%%%%%%%%%%%%%%%%%%%%%%%%%%%%%%%%%%%%%%%%%%%%%%%%%%%%%%%%%%%%%%%%%%%%

\textwidth=31.9pc
\textheight=46.5pc
\oddsidemargin=1pc
\evensidemargin=1pc
\headsep=15pt
\topmargin=.6cm
\parindent=1.7pc
\parskip=0pt

\usepackage{amsmath}
\usepackage{amssymb}
\usepackage{amsfonts}
\usepackage{multirow}
\usepackage{amsthm}

% Marcacao de alteracoes a partir de k = 2 (docs/instrucoes.md, S3): a cor da
% rodada corrente e colR1, e o ulem fornece o \sout das remocoes sugeridas.
% A versao k = 1 nao leva marcacao, porque nao ha versao anterior.
\usepackage{xcolor}
\usepackage[normalem]{ulem}
\definecolor{colR1}{RGB}{0,90,200}

\setcounter{page}{1}
\newtheorem{theorem}{Theorem}
\newtheorem{lemma}{Lemma}
\newtheorem{corollary}{Corollary}
\newtheorem{proposition}{Proposition}
\theoremstyle{definition}
\newtheorem{definition}{Definition}
\newtheorem{example}{Example}
\newtheorem{remark}{Remark}
\newtheorem{assumption}{Assumption}
\pagestyle{fancy}

% --- macros -------------------------------------------------------------
% Implementam docs/notacao.md, como derivations/macros.tex, com tres
% diferencas exigidas pelo estilo da revista (S5 de ss-instrucoes-autores.md):
% E(.) e P(.) em romano e com parenteses, no lugar do \mathbb{E}[.] das
% derivacoes, e B_X no lugar dos dois simbolos C_X (E1.3) e B_X (E1.4) para a
% cota das covariaveis lineares.
\newcommand{\R}{\mathbb{R}}
\newcommand{\E}{\mathrm{E}}
\newcommand{\Prob}{\mathrm{P}}
\newcommand{\Var}{\mathrm{Var}}
\newcommand{\T}{^{\top}}
\newcommand{\norm}[1]{\left\lVert #1 \right\rVert}
\newcommand{\normn}[1]{\left\lVert #1 \right\rVert_{n}}
\newcommand{\normL}[1]{\left\lVert #1 \right\rVert_{L_2(P)}}
\newcommand{\inner}[2]{\langle #1, #2 \rangle}
\newcommand{\cv}{\ell}
\newcommand{\md}{m}
\newcommand{\lev}{j}
\newcommand{\tsl}{k}
\newcommand{\jzero}{j_{0}}
\newcommand{\NJ}{N_{J}}
\newcommand{\bX}{\mathbf{X}}
\newcommand{\bU}{\mathbf{U}}
\newcommand{\bZ}{\mathbf{Z}}
\newcommand{\fcoef}[1]{\beta_{#1}}
\newcommand{\lvl}[1]{c_{#1}}
\newcommand{\gcomp}[2]{g_{#1#2}}
\newcommand{\wav}[2]{\psi_{#1#2}}
\newcommand{\VJ}{V_{J}}
\newcommand{\WJ}{W_{J}}
\newcommand{\PiJ}{\Pi_{J}}
\newcommand{\coef}[4]{\theta_{#1#2,#3#4}}
\newcommand{\coefs}[4]{\theta^{*}_{#1#2,#3#4}}
\newcommand{\btheta}{\boldsymbol{\theta}}
\newcommand{\thetastar}{\btheta^{*}}
\newcommand{\thetahat}{\widehat{\btheta}}
\newcommand{\bbeta}{\boldsymbol{\beta}}
\newcommand{\betastar}{\bbeta^{*}}
\newcommand{\bc}{\mathbf{c}}
\newcommand{\chat}{\widehat{\mathbf{c}}}
\newcommand{\Gram}{\Sigma}
\newcommand{\Gramhat}{\widehat{\Sigma}}
\newcommand{\pen}{\lambda}
\newcommand{\penn}{\lambda_{n}}
\newcommand{\Szero}{S_{0}}
\newcommand{\Besov}[3]{B^{#1}_{#2,#3}}
\newcommand{\seff}{s'}
\newcommand{\Amat}{\mathbf{A}}
\newcommand{\Bmat}{\mathbf{B}}
\newcommand{\Btil}{\widetilde{\Bmat}}
\newcommand{\PA}{\Pi_{\Amat}}
\newcommand{\MA}{M_{\Amat}}
\newcommand{\bveps}{\boldsymbol{\varepsilon}}
\newcommand{\bb}{\mathbf{b}}
\newcommand{\bY}{\mathbf{Y}}
\newcommand{\bPsi}{\boldsymbol{\Psi}}
\newcommand{\smax}{\widehat{\sigma}_{\max}}
\newcommand{\gtil}{\widetilde{\gamma}}
\newcommand{\Mstar}{M^{\star}}
\newcommand{\lmin}{\lambda_{\min}}
\newcommand{\lmaxx}{\lambda_{\max}}
\newcommand{\opnorm}[1]{\left\lVert #1 \right\rVert_{\mathrm{op}}}

%%%%%%%%%%%%%%%%%%%%%%%%%%%%%%%%%%%%%%%%%%%%%%%%%%%%%%%%%%%%%%%%%%%%%%%%%%%%%%
\pagestyle{fancy}
\def\n{\noindent}
\lhead[\fancyplain{} \leftmark]{}
\chead[]{}
\rhead[]{\fancyplain{}\rightmark}
\cfoot{}

%%%%%%%%%%%%%%%%%%%%%%%%%%%%%%%%%%%%%%%%%%%%%%%%%%%%%%%%%%%%%%%%%%%%%%%%%%%%%%
\begin{document}

\renewcommand{\baselinestretch}{2}

\markright{ \hbox{\footnotesize\rm Statistica Sinica
  }\hfill\\[-13pt]
  \hbox{\footnotesize\rm
  }\hfill }

\markboth{\hfill{\footnotesize\rm AUTHOR(S)} \hfill}
{\hfill {\footnotesize\rm WAVELET ADDITIVE FUNCTIONAL COEFFICIENTS} \hfill}

\renewcommand{\thefootnote}{}
$\ $\par

%%%%%%%%%%%%%%%%%%%%%%%%%%%%%%%%%%%%%%%%%%%%%%%%%%%%%%%%%%%%%%%%%%%%%%%%%%%%%%

\fontsize{12}{14pt plus.8pt minus .6pt}\selectfont \vspace{0.8pc}
\centerline{\large\bf SPARSE WAVELET ESTIMATION OF}
\vspace{2pt}
\centerline{\large\bf ADDITIVE FUNCTIONAL COEFFICIENTS}
\vspace{.4cm}
\centerline{Author(s)}
\vspace{.4cm}
\centerline{\it Affiliation(s)}
\vspace{.55cm} \fontsize{9}{11.5pt plus.8pt minus.6pt}\selectfont

%%%%%%%%%%%%%%%%%%%%%%%%%%%%%%%%%%%%%%%%%%%%%%%%%%%%%%%%%%%%%%%%%%%%%%%%%%%%%%

\begin{quotation}
\noindent {\it Abstract:}
We study a varying coefficient model in which the effect of each linear
covariate is additive in several modulating covariates, and each additive
component is expanded in a periodized wavelet basis. The resulting design is
a sieve of products between a linear covariate and a wavelet evaluated at a
modulator, and the components are estimated by a single lasso problem in
which the levels of the coefficients are left unpenalized. Three results are
established. The population Gram matrix of the product design has its
smallest eigenvalue bounded away from zero under a conditional second moment
condition on the linear covariates and a bounded joint density for the
modulators, with no independence between the two groups of covariates; the
estimator satisfies a nonasymptotic oracle inequality that needs no cone
condition, because the unpenalized levels are profiled out; and, at a
resolution level chosen inside an explicit window, the prediction error and
the componentwise error attain the rate of nonlinear wavelet approximation up
to a logarithmic factor, under the same Besov condition that governs the
linear sieve. The last result rests on the observation that the Besov
condition already implies compressibility of the sieve coefficients, so that
adaptivity costs no additional assumption. The estimator is a single convex
program and is fit by coordinate descent on a sparse design matrix.

\vspace{9pt}
\noindent {\it Key words and phrases:}
Additive coefficient model, Besov space, effective dimension, lasso, oracle
inequality, varying coefficient model, wavelets.
\par
\end{quotation}\par

\def\thefigure{\arabic{figure}}
\def\thetable{\arabic{table}}

\renewcommand{\theequation}{\thesection.\arabic{equation}}
% o contador de equacoes do template nao reinicia por secao, o que produzia
% "(3.6)" logo depois de "(2.5)"; o \numberwithin reinicia.
\numberwithin{equation}{section}

\fontsize{12}{14pt plus.8pt minus .6pt}\selectfont

%%%%%%%%%%%%%%%%%%%%%%%%%%%%%%%%%%%%%%%%%%%%%%%%%%%%%%%%%%%%%%%%%%%%%%%%%%%%%%
\section{Introduction}
\label{sec:intro}

Applied problems often present an effect that is not constant across the
sample. The return to schooling changes with age and with experience, the
effect of temperature on electricity demand changes with the hour of the day
and with humidity, and the effect of a treatment changes with the age and the
body mass index of the patient. The varying coefficient model of
\citet{Hastie-Tibshirani-1993} and \citet*{Cai-Fan-Yao-2000} accommodates this
by letting the coefficient of each covariate be a function of a modulating
variable; see \citet{Fan-Zhang-2008} for a review. When more than one variable
modulates the effect, a fully nonparametric coefficient runs into the curse of
dimensionality, and the standard remedy is to make each coefficient additive
in the modulators, which is the additive coefficient model of
\citet{Xue-Yang-2006}, studied further by \citet{Xue-Liang-2010} and
\citet{Liu-Yang-2010}.

The estimators available for this model are built from polynomial splines and
from kernels, and are governed by a global smoothness index, so that a
component with a peak, a jump or a change of regime is either oversmoothed or
paid for everywhere in the domain. Wavelet bases combined with an $\ell_1$
penalty are the standard answer to this difficulty, because a function whose
smoothness is inhomogeneous is still economically represented by a few large
coefficients \citep{Donoho-Johnstone-1994, Donoho-Johnstone-1998,
Antoniadis-Fan-2001}. They have been used for varying coefficients without an
additive structure and without penalization \citep*{Zhou-You-2004,
Montoril-Morettin-Chiann-2018}, and for additive models without a multiplying
covariate \citep{Sardy-Tseng-2004, Zhang-Wong-2003, Haris-Simon-Shojaie-2018,
Amato-Antoniadis-DeFeis-Gijbels-2022, Sardy-Ma-2024}; penalized estimation in
the additive coefficient model itself has been carried out with grouped
penalties on spline bases \citep{Wei-Huang-Li-2011,
Antoniadis-Gijbels-LambertLacroix-2014, Ma-Carroll-Liang-Xu-2015}, again under
global smoothness.

\citet{Klopp-Pensky-2015} study the varying coefficient model in which a
single index modulates every coefficient. They expand each coefficient
function in an orthonormal basis, estimate the resulting array by a block
lasso, and obtain a nonasymptotic oracle inequality, a Besov rate that adapts
to inhomogeneous smoothness, and a matching minimax lower bound, under the
assumption that the covariates are independent of the index. The present paper
carries that programme to the situation an applied problem usually presents,
in which several variables modulate a coefficient at once. We let each
coefficient be additive in the modulators, which keeps the estimator free of
the curse of dimensionality in their number, we allow the covariates to depend
on the modulators, and we penalize coefficient by coefficient rather than in
blocks, leaving the levels of the coefficients unpenalized. Each of these
costs something. Additivity produces cross terms between distinct modulators,
so the population Gram matrix no longer factorizes as a Kronecker product and
has to be bounded directly; dependence between the covariates and the
modulators replaces a marginal second moment by a conditional one; and the
unpenalized levels leave a term in the risk that block penalization does not
incur. Our main result, Theorem~\ref{thm:main}, states that the estimator
attains the rate of nonlinear wavelet approximation up to a logarithmic
factor, under the same Besov assumption that governs the linear sieve, and it
reduces to the rate of Klopp and Pensky when a single modulator is present.

We call the resulting procedure WAFC, for wavelet additive functional
coefficients. Its contributions are of three kinds. The first is the theory of
the additive product design: a design condition for the Gram matrix of the
products between a linear covariate and a wavelet of a modulator, which
carries a cross term between distinct modulators and does not assume
independence (Proposition~\ref{prop:gram}); an oracle inequality that needs no
cone condition, because the unpenalized levels are profiled out
(Theorem~\ref{thm:oracle}); and rates obtained by compressibility
(Theorem~\ref{thm:main}). The second is the reading that compressibility costs
no assumption: the Besov condition on the components already implies that
their wavelet coefficients lie in a weak $\ell_\tau$ ball with
$\tau = (s+1/2)^{-1}$ (Lemma~\ref{lem:weak}), and it is this that separates
the rate of the estimator from the rate of the linear sieve when the Besov
integrability index is smaller than two. The third is computational: the
estimator is one convex program on a sparse design matrix, and the resolution
level and the penalty are chosen jointly.

The rest of the paper is organized as follows. Section~\ref{sec:model}
introduces the model, its identification, the wavelet sieve and the estimator.
Section~\ref{sec:theory} states the assumptions and the theory, opening with
the main rate and then giving the three results it rests on.
Section~\ref{sec:computation} describes the computation, the scale of the
penalty and the tuning rules. Proofs are collected in the online Supplementary
Material.
% E5b: acrescentar aqui a mencao as Secoes 5 (Simulation), 6 (Application) e
% 7 (Discussion), e a frase de fechamento do roteiro.

Throughout, $\norm{a}_\pi$ is the $\ell_\pi$ norm of a vector,
$\normn{h}^2 = n^{-1}\sum_{i\le n} h_i^2$ is the empirical norm,
$\lmin(\cdot)$ and $\lmaxx(\cdot)$ are the extreme eigenvalues of a symmetric
matrix, $\opnorm{\cdot}$ is the operator norm, and $a_n \asymp b_n$ means that
$a_n/b_n$ is bounded away from zero and infinity. For a support
$S$, $\delta_S$ is the vector that agrees with $\delta$ on $S$ and vanishes
elsewhere.

%%%%%%%%%%%%%%%%%%%%%%%%%%%%%%%%%%%%%%%%%%%%%%%%%%%%%%%%%%%%%%%%%%%%%%%%%%%%%%
\section{Model and Wavelet Sieve}
\label{sec:model}

\subsection{The additive functional coefficient model}
\label{sec:model-def}

Let $Y$ be a scalar response, $\bX = (X_1,\dots,X_p)\T$ a vector of linear
covariates and $\bU = (U_1,\dots,U_q)\T$ a vector of modulating covariates,
taking values in $[0,1]^q$ after a monotone transformation of each
coordinate. The model is
\begin{equation}
  \label{eq:model}
  Y \;=\; \sum_{\cv=1}^{p} \fcoef{\cv}(\bU)\,X_{\cv} \;+\; \varepsilon,
  \qquad \E(\varepsilon \mid \bX, \bU) = 0,
\end{equation}
with each functional coefficient additive in the modulators,
\begin{equation}
  \label{eq:additive}
  \fcoef{\cv}(u) \;=\; \lvl{\cv} \;+\; \sum_{\md=1}^{q}\gcomp{\cv}{\md}(u_{\md}),
  \qquad \int_{0}^{1}\gcomp{\cv}{\md}(v)\,dv = 0,
  \qquad \cv = 1,\dots,p .
\end{equation}
The constant $\lvl{\cv}$ is the level of the $\cv$th coefficient and
$\gcomp{\cv}{\md}$ is the component of that coefficient along the $\md$th
modulator. The centering in \eqref{eq:additive} is with respect to Lebesgue
measure, which is the convention of the wavelet literature and the one the
basis of Section~\ref{sec:sieve} imposes at no cost; the parametrization
centered at the law of $\bU$ is obtained from it by the shift
$\lvl{\cv} \mapsto \lvl{\cv} + \sum_{\md}\E\{\gcomp{\cv}{\md}(U_{\md})\}$,
which moves mass between the levels and the components but changes neither
$\fcoef{\cv}$ nor the estimator.

The model contains several familiar ones. With $X_1 \equiv 1$ the first
coefficient is an additive intercept, so that $p = 1$ gives the additive model
and $p > 1$ with $\gcomp{\cv}{\md} \equiv 0$ for $\cv \ge 2$ gives the
partially linear additive model; $q = 1$ gives the classical varying
coefficient model with a single index; and $\gcomp{\cv}{\md}\equiv 0$ for all
pairs gives linear regression. Writing
$f(x,u) = \sum_{\cv} x_{\cv}\fcoef{\cv}(u)$ for the regression function, the
first question is whether the distribution of $(Y,\bX,\bU)$ determines the
parameters. Two degeneracies have to be excluded: if $\bX$ is confined to a
hyperplane given $\bU$, only a combination of the coefficients is determined,
and if one modulator is a function of another, only the sum of the
corresponding components is.

% derivations/01-identificabilidade.md, Proposicao 1 (numeracao global).
\begin{proposition}[Identification]
\label{prop:identif}
Suppose that $\E\norm{\bX}_2^2 < \infty$, that $\E(\bX\bX\T \mid \bU)$ is
nonsingular almost surely, and that $\bU$ has a Lebesgue density on $[0,1]^q$
that is positive almost everywhere. Then, among the parametrizations that
satisfy \eqref{eq:additive}, the levels $\lvl{1},\dots,\lvl{p}$ and the
components $\gcomp{\cv}{\md} \in L_2[0,1]$ are determined by $f$: the levels
uniquely, and each component up to a Lebesgue null set.
\end{proposition}

Dependence between $\bX$ and $\bU$ is allowed, and so is dependence among the
modulators; what the two conditions exclude is a linear combination of the
linear covariates that is a function of the modulators, and a modulator that
is a function of the others. Both exclusions return in quantitative form in
Section~\ref{sec:theory}.

\subsection{The wavelet sieve}
\label{sec:sieve}

Let $\phi$ and $\psi$ be the scaling function and the mother wavelet of an
orthonormal multiresolution analysis of $L_2(\R)$, compactly supported, with
$\psi$ bounded and with $N$ vanishing moments, and let $L$ be the length of
the filter, so that the support of $\psi$ is contained in $[0,L-1]$. Write
\begin{equation}
  \label{eq:periodized}
  \wav{\lev}{\tsl}(u) \;=\; \sum_{\iota \in \mathbb{Z}}
  2^{\lev/2}\psi\{2^{\lev}(u+\iota) - \tsl\},
  \qquad \lev \ge 0, \quad \tsl = 0,\dots,2^{\lev}-1,
\end{equation}
for the periodization of the basis to $[0,1]$. The system
$\{\mathbf{1}\}\cup\{\wav{\lev}{\tsl}\}$ is an orthonormal basis of
$L_2[0,1]$ under Lebesgue measure and every wavelet in it integrates to zero
\citep[Section 9.3]{Daubechies-1992}. We take the coarsest level to be
$\jzero = 0$, so that the only scaling function is the constant
$\phi_{00}\equiv 1$, and we discard it: it is already carried by the level
$\lvl{\cv}$. The sieve of resolution $J$ is then
\[
  \WJ \;=\; \mathrm{span}\{\wav{\lev}{\tsl} : 0 \le \lev < J,\;
  0 \le \tsl < 2^{\lev}\},
  \qquad \dim \WJ = \NJ = 2^{J}-1,
\]
the orthogonal complement of the constants inside the multiresolution space
$\VJ$, and $\PiJ$ is the orthogonal projection onto it. Every element of $\WJ$
integrates to zero, so that the identification constraint of
\eqref{eq:additive} holds by construction, no numerical restriction is imposed
on the fit, and the sieve is a plain linear space.

% derivations/01-identificabilidade.md, Lema 1 (numeracao global).
\begin{lemma}[The basis carries the constraint]
\label{lem:basis}
Under the conditions of Proposition~\ref{prop:identif}, for every $J$ the map
\[
  (\bc, \btheta) \;\longmapsto\;
  \sum_{\cv=1}^{p} x_{\cv}\Bigl\{\lvl{\cv} + \sum_{\md=1}^{q}
  \sum_{\lev<J}\sum_{\tsl} \coef{\cv}{\md}{\lev}{\tsl}
  \wav{\lev}{\tsl}(u_{\md})\Bigr\}
\]
is injective on $\R^{p}\times\R^{d}$, $d = pq\NJ$; equivalently, the
population Gram matrix of the design defined in \eqref{eq:design} is positive
definite. If instead the column $X_{\cv}\phi_{00}(U_{\md})$ is retained in
each block, that matrix is singular with nullity exactly $pq$.
\end{lemma}

Replacing $\gcomp{\cv}{\md}$ by $\PiJ\gcomp{\cv}{\md}$ leaves a bias, which
Section~\ref{sec:theory} controls under a Besov condition. Two practical
consequences of periodization deserve mention. A component whose values at the
two ends of the interval disagree is not periodic, and its projection error
then stalls at the order $2^{-J/2}$ in $L_2$ regardless of the interior
smoothness (Proposition~S3.1 in the Supplementary Material). The two remedies,
rescaling the modulators to a proper subinterval and using the basis of
\citet*{cohen1993wavelets} adapted to the interval, are matters of computation
and are discussed in Section~\ref{sec:computation}.

\subsection{The estimator}
\label{sec:estimator}

Let $(Y_i,\bX_i,\bU_i)$, $i = 1,\dots,n$, be independent copies of
$(Y,\bX,\bU)$. The row of the design matrix is the vector of the $p$ linear
covariates followed by the $d = pq\NJ$ products between a linear covariate and
a wavelet of a modulator,
\begin{equation}
  \label{eq:design}
  \bZ = [\,\Amat \;\; \Bmat\,], \qquad
  \Amat_{i\cv} = X_{i\cv}, \qquad
  \Bmat_{i,(\cv\md,\lev\tsl)} = X_{i\cv}\,\wav{\lev}{\tsl}(U_{i\md}),
\end{equation}
with the blocks $(\cv,\md)$ in lexicographic order and, inside a block, the
wavelets ordered by level and then by translation. Writing
$\bY = (Y_1,\dots,Y_n)\T$, the estimator is
\begin{equation}
  \label{eq:estimator}
  (\chat,\thetahat) \;\in\; \mathop{\rm arg\,min}_{(\bc,\btheta)\in
  \R^{p}\times\R^{d}}
  \Bigl\{\normn{\bY - \Amat\bc - \Bmat\btheta}^{2}
  \;+\; 2\pen\norm{\btheta}_{1}\Bigr\} .
\end{equation}
The penalty does not reach the levels $\bc$: they are $p$ parameters of a
parametric part, and shrinking them would bias every fitted coefficient by a
constant. The fitted components and coefficients are recovered from
$\thetahat$ by
\[
  \widehat{g}_{\cv\md}(v) = \sum_{\lev<J}\sum_{\tsl}
  \widehat{\theta}_{\cv\md,\lev\tsl}\,\wav{\lev}{\tsl}(v),
  \qquad
  \widehat{\fcoef{\cv}}(u) = \widehat{c}_{\cv}
  + \sum_{\md=1}^{q}\widehat{g}_{\cv\md}(u_{\md}).
\]
Problem \eqref{eq:estimator} is a lasso with some coordinates unpenalized, so
it is convex, and it is strictly convex, hence uniquely solved, whenever the
Gram matrix of $\bZ$ is nonsingular. Two structural variants are available at
no conceptual cost and are discussed in Section~\ref{sec:computation}: a
sparse group lasso \citep{Simon-Friedman-Hastie-Tibshirani-2013} whose groups
are the blocks $(\cv,\md)$, which removes an entire component when a modulator
does not act on a coefficient, and the block lasso of
\citet{Klopp-Pensky-2015}. The theory below is for \eqref{eq:estimator}.

%%%%%%%%%%%%%%%%%%%%%%%%%%%%%%%%%%%%%%%%%%%%%%%%%%%%%%%%%%%%%%%%%%%%%%%%%%%%%%
\section{Theory}
\label{sec:theory}

\subsection{Assumptions}
\label{sec:assumptions}

\begin{assumption}[Sampling and error]
\label{ass:sample}
The observations $(Y_i,\bX_i,\bU_i)$, $i \le n$, are independent and
identically distributed and follow \eqref{eq:model}, and the error is
conditionally sub-Gaussian: there is $\sigma < \infty$ with
$\E\{\exp(t\varepsilon_i) \mid \bX_i,\bU_i\} \le \exp(\sigma^2t^2/2)$ almost
surely, for all real $t$.
\end{assumption}

\begin{assumption}[Linear covariates]
\label{ass:X}
$\max_{\cv}|X_{\cv}| \le B_X$ almost surely, and there are constants
$0 < \kappa_1 \le \kappa_2 < \infty$ with
$\kappa_1 \le \lmin\{\E(\bX\bX\T\mid\bU)\} \le
\lmaxx\{\E(\bX\bX\T\mid\bU)\} \le \kappa_2$ almost surely.
\end{assumption}

\begin{assumption}[Modulating covariates]
\label{ass:U}
$\bU$ takes values in $[0,1]^q$ and has a joint Lebesgue density $p_{\bU}$
with $0 < c_U \le p_{\bU}(u) \le C_U < \infty$ for every $u \in [0,1]^q$.
\end{assumption}

\begin{assumption}[Basis]
\label{ass:basis}
The wavelets are the periodization \eqref{eq:periodized} of an orthonormal,
compactly supported family with $\psi$ bounded and $N$ vanishing moments, and
$\jzero = 0$.
\end{assumption}

\begin{assumption}[Besov regularity]
\label{ass:besov}
There are $s > 0$, $1 \le \pi, r \le \infty$ and $C_g < \infty$ such that
every component satisfies $\int_0^1 \gcomp{\cv}{\md} = 0$ and, writing
$\coefs{\cv}{\md}{\lev}{\tsl} = \inner{\gcomp{\cv}{\md}}{\wav{\lev}{\tsl}}$
and $\theta^{*}_{\cv\md,\lev\cdot}$ for the vector of the coefficients of
level $\lev$,
\[
  \norm{\theta^{*}_{\cv\md,\lev\cdot}}_{\pi} \;\le\;
  C_g\,2^{-\lev(s + 1/2 - 1/\pi)} \qquad \text{for every } \lev \ge 0 .
\]
The effective regularity $\seff = s - (1/\pi - 1/2)_{+}$ is positive.
\end{assumption}

Assumptions~\ref{ass:X} and \ref{ass:U} are the quantitative forms of the two
conditions of Proposition~\ref{prop:identif}. In the second it is the joint
density that is bounded below, not the marginals: if $U_2 = U_1$ almost surely
both marginals are uniform and the design is singular.
Assumption~\ref{ass:besov} is taken as a condition on the sequence of
coefficients rather than as membership in a function space, as in
\citet{Klopp-Pensky-2015}; when $\psi$ has more than $s$ vanishing moments and
enough Hölder regularity the two are equivalent, on the circle
\citep[Section 3.6]{Cohen-2003}, and the implication that the proofs use needs
only the vanishing moments. The index $\seff$ is the regularity that a linear
sieve sees and $s$ is the one that nonlinear approximation sees; they differ
exactly when $\pi < 2$, which is the regime of inhomogeneous smoothness.

\subsection{The main result}
\label{sec:main}

Write $\widehat{f}_i = \Amat_i\chat + \Bmat_i\thetahat$ for the fit, $d_n$ for
the number of penalized columns at resolution $J_n$, and fix a confidence
level $\alpha \in (0,1)$. The penalty is the one that the oracle inequality of
Section~\ref{sec:oracle} calls for, namely
\begin{equation}
  \label{eq:lambda}
  \penn \;=\; 2\sigma B_X\sqrt{2C_U}\,
  \sqrt{\frac{2\log(2d_n/\alpha)}{n}} \;\asymp\;
  \sigma B_X \sqrt{C_U}\,\sqrt{\frac{J_n}{n}} .
\end{equation}

% derivations/05-taxas.tex, Corolario 5 (numeracao global); D16.
\begin{theorem}[Rate of nonlinear approximation]
\label{thm:main}
Let Assumptions~\ref{ass:sample} to \ref{ass:besov} hold with $p$, $q$ and all
the constants fixed as $n$ grows, and set $\tau = (s+1/2)^{-1}$. Suppose
$\seff > s/(2s+1)$ and take $J_n = \lceil c \log_2 n\rceil$ with
\begin{equation}
  \label{eq:window}
  \max\Bigl\{\frac{\tau}{2},\; \frac{2-\tau}{4\seff}\Bigr\}
  \;\le\; c \;<\; 1 ,
\end{equation}
a window that is nonempty precisely under that condition. Then, with $\penn$
as in \eqref{eq:lambda},
\[
  \normn{\widehat{f}-f}^{2}
  \;=\; O_p\Bigl\{\Bigl(\frac{\log n}{n}\Bigr)^{2s/(2s+1)}\Bigr\},
  \qquad
  \sum_{\cv,\md}\bigl\lVert \widehat{g}_{\cv\md}-\gcomp{\cv}{\md}
  \bigr\rVert_{L_2[0,1]}^{2}
  \;=\; O_p\Bigl\{\Bigl(\frac{\log n}{n}\Bigr)^{2s/(2s+1)}\Bigr\} .
\]
\end{theorem}

Three readings of Theorem~\ref{thm:main} matter. First, the exponent is the
one of nonlinear wavelet approximation, $2s/(2s+1)$, and not the one of the
linear sieve, $2\seff/(2\seff+1)$ of Theorem~\ref{thm:sieve} below; the two
agree when $\pi \ge 2$ and separate when $\pi < 2$, which is where the
$\ell_1$ penalty pays for itself. Second, the rate is attained under
Assumption~\ref{ass:besov} alone: no compressibility is assumed on top of it,
because the Besov condition already implies it, with
$\tau = (s+1/2)^{-1}$, by Lemma~\ref{lem:weak}. Third, the rate
$(\log n/n)^{2s/(2s+1)}$ is the minimax reference of the Gaussian sequence
model over weak $\ell_\tau$ balls with per coordinate noise of order $\penn$
\citep{Donoho-Johnstone-1998}; the logarithmic factor is the price of the
single, nonadaptive penalty level \eqref{eq:lambda}, which takes a union bound
over the $d_n$ columns. No lower bound is proved here for $q \ge 2$: the
comparison is with the reference of the sequence model, and the minimax lower
bound of \citet{Klopp-Pensky-2015} covers the case of a single modulator with
covariates independent of it.

The proof combines three ingredients, given below: the approximation error of
the sieve, the design condition, and a nonasymptotic oracle inequality.
Section~\ref{sec:rates} assembles them, first at a dense reading of the
estimation term, which gives the sieve rate, and then at a sparse comparator,
which gives Theorem~\ref{thm:main}.

\subsection{Approximation and the design condition}
\label{sec:approx}

% derivations/02-aproximacao-besov.tex, Lema 2 e Corolario 1 (numeracao global).
\begin{lemma}[Bias of the sieve]
\label{lem:bias}
Under Assumptions~\ref{ass:basis} and \ref{ass:besov}, for every pair
$(\cv,\md)$ and every $J \ge 0$,
\[
  \bigl\lVert \gcomp{\cv}{\md} - \PiJ\gcomp{\cv}{\md}
  \bigr\rVert_{L_2[0,1]}^{2}
  \;=\; \sum_{\lev \ge J}\norm{\theta^{*}_{\cv\md,\lev\cdot}}_{2}^{2}
  \;\le\; \frac{C_g^{2}}{1-2^{-2\seff}}\,2^{-2J\seff} .
\]
If in addition Assumptions~\ref{ass:X} and \ref{ass:U} hold, then, with
$f_J(x,u) = \sum_{\cv}x_{\cv}\{\lvl{\cv} +
\sum_{\md}(\PiJ\gcomp{\cv}{\md})(u_{\md})\}$,
\begin{equation}
  \label{eq:bias}
  \bigl\lVert f - f_J\bigr\rVert_{L_2(P)}^{2} \;\le\;
  \frac{(pq)^{2}B_X^{2}C_UC_g^{2}}{1-2^{-2\seff}}\,2^{-2J\seff}
  \;=:\; \mathcal{B}_{J}^{2} .
\end{equation}
\end{lemma}

The exponent is $\seff$ and not $s$, and this is where the loss
$(1/\pi-1/2)_{+}$ enters: inside a level, an $\ell_\pi$ bound with $\pi < 2$
controls the $\ell_2$ norm without gain. The bias uses only the bound on the
linear covariates and the marginal density of each modulator, so dependence
between $\bX$ and $\bU$ does not affect it. What it does affect is the passage
from the prediction norm to the norm of each component, which is the design
condition.

Let $\bPsi(u) \in \R^{K}$, $K = 1 + q\NJ$, collect the constant and the
wavelets of the $q$ modulators, so that the row of $\bZ$ in
\eqref{eq:design} is a permutation of the Kronecker product
$\bX \otimes \bPsi(\bU)$. Write $\Gram = \E(\bZ\bZ\T)$ and
$\Gramhat = \bZ\T\bZ/n$ for the population and empirical Gram matrices and
$\Gram_{\bPsi} = \E\{\bPsi(\bU)\bPsi(\bU)\T\}$ for the part in the modulators
alone.

% derivations/03-desenho-produtos.tex, Proposicao 3 (numeracao global).
\begin{proposition}[Gram matrix of the product design]
\label{prop:gram}
Let Assumptions~\ref{ass:X} to \ref{ass:basis} hold and write $D = p + d$ and
$\gamma = \kappa_1 c_U$.
\begin{enumerate}
\item[(i)] $c_U \le \lmin(\Gram_{\bPsi}) \le \lmaxx(\Gram_{\bPsi}) \le C_U$
for every $J$, and consequently
$\gamma \le \lmin(\Gram) \le \lmaxx(\Gram) \le \kappa_2 C_U$. If $\bX$ and
$\bU$ are independent, $\Gram$ is, up to the permutation of columns, the
Kronecker product $\E(\bX\bX\T)\otimes\Gram_{\bPsi}$, and the two bounds are
attained as products of eigenvalues.
\item[(ii)] With $R_J = pB_X^{2}(1 + qC_\psi 2^{J})$, where $C_\psi$ depends
only on the wavelet family, $\norm{\bZ_i}_2^{2} \le R_J$ almost surely and
\[
  \Prob\Bigl(\lmin(\Gramhat) \ge \frac{\gamma}{2}\Bigr)
  \;\ge\; 1 - 2D\exp\Bigl\{-\frac{n\gamma^{2}}
  {8R_J(\kappa_2C_U + \gamma/3)}\Bigr\},
\]
which tends to one whenever $pq2^{J}\log(pq2^{J})/n \to 0$.
\end{enumerate}
\end{proposition}

Part (i) is stated for the general case: independence between the covariates
and the modulators buys the exact factorization but not a better constant. A
lower bound on the smallest eigenvalue dominates every compatibility and
restricted eigenvalue constant, on every support and with no cone restriction,
since $\delta\T A \delta \ge \lmin(A)\norm{\delta}_2^2 \ge
\lmin(A)\norm{\delta_S}_1^2/|S|$, and it is in this form that the next
subsection consumes it. In the resolution regime of Section~\ref{sec:rates}
the requirement of part (ii) is of the same order as the rate itself, so it
costs nothing there; if $pq$ were larger than $n^{1/2}$, only the
compatibility constant could be transferred, through the entrywise deviation
\citep[Corollary 6.8]{buhlmann2011statistics}, at the price of a factor
$s_0^{2}$ and of the cone.

\subsection{A nonasymptotic oracle inequality}
\label{sec:oracle}

Let $\betastar = (\bc,\thetastar)$ be the sieve oracle, that is the vector
whose components are the levels and the coefficients
$\coefs{\cv}{\md}{\lev}{\tsl}$ of Assumption~\ref{ass:besov} for $\lev < J$,
so that $f_J$ is its fit, and let
\[
  \bb = f - f_J \in \R^{n}, \qquad \Szero = \mathrm{supp}(\thetastar),
  \qquad s_0 = |\Szero|, \qquad v = \thetahat - \thetastar .
\]
Let $\PA$ be the orthogonal projection onto the column space of $\Amat$, which
exists whenever $\Amat$ has rank $p$, let $\MA = I_n - \PA$, and let
$\Btil = \MA\Bmat$ and $\gtil = \lmin(\Btil\T\Btil/n)$ be the residualized
design and its smallest eigenvalue. Finally let
$\smax^{2} = \max_a (\Gramhat)_{aa}$ be the largest empirical norm of a
penalized column.

% derivations/04-oraculo.tex, Teorema 1 (numeracao global).
\begin{theorem}[Oracle inequality]
\label{thm:oracle}
Let Assumption~\ref{ass:sample} hold, suppose $\Amat$ has rank $p$, and take
$\pen \ge 2\pen_0$ with
$\pen_0 = \sigma\smax\{2\log(2d/\alpha)/n\}^{1/2}$. On the event
$\mathcal{T} = \{\norm{\Btil\T\bveps/n}_{\infty} \le \pen/2\}$, which has
probability at least $1 - \alpha$:
\begin{enumerate}
\item[(i)] with no condition on the design,
$\normn{\Btil v}^{2} \le 6\pen\norm{\thetastar}_{1} + 4\normn{\bb}^{2}$;
\item[(ii)] if $\gtil > 0$,
\[
  \normn{\Btil v}^{2} \le \frac{64\pen^{2}s_0}{\gtil} + 16\normn{\bb}^{2},
  \qquad
  \norm{v}_{2}^{2} \le \frac{64\pen^{2}s_0}{\gtil^{2}}
  + \frac{16\normn{\bb}^{2}}{\gtil} ;
\]
\item[(iii)] in either regime,
$\normn{\widehat{f}-f} \le \normn{\Btil v} + 2\normn{\bb}
+ \normn{\PA\bveps}$, and
$\E(\normn{\PA\bveps}^{2}\mid\bX,\bU) \le \sigma^{2}p/n$.
\end{enumerate}
\end{theorem}

Three features separate Theorem~\ref{thm:oracle} from the standard treatment
\citep[Chapter 6]{buhlmann2011statistics}. The unpenalized levels are not
placed inside the support, which would require a compatibility condition on a
cone; they are profiled out, and the Schur complement transfers the full
smallest eigenvalue of $\Gramhat$ to the residualized block,
$\gtil \ge \lmin(\Gramhat)$, so that Proposition~\ref{prop:gram}(ii) is used
as it stands and the constant does not depend on $\Szero$. The model is
misspecified inside the sieve, and the deterministic residual $\bb$ is handled
by Young's inequality rather than by duality, which keeps the statement
nonasymptotic and valid also when the bias dominates. The penalty is
calibrated by the empirical column norms and not by $\max_{i,a}|Z_{ia}|$: the
former is bounded in probability by $2B_X^2C_U$, while the latter is of order
$2^{J/2}$ and would inflate $\pen$ by that factor and destroy the rates of
Section~\ref{sec:rates}. The term $\sigma^{2}p/n$ in part (iii) is the price
of leaving the $p$ levels out of the penalty; no choice of $\pen$ removes it,
it is of smaller order than the estimation term whenever
$p \lesssim s_0\log(pq\NJ)$, and it is absent from
\citet{Klopp-Pensky-2015}, where the levels are penalized along with
everything else.

\subsection{Rates}
\label{sec:rates}

Under Assumption~\ref{ass:besov} the sieve oracle is generically dense: the
coefficients decay but do not vanish, so that $s_0 = d$ and part (ii) of
Theorem~\ref{thm:oracle}, read at the oracle, delivers the rate of the linear
sieve.

% derivations/05-taxas.tex, Teorema 2 (numeracao global).
\begin{theorem}[Rate of the linear sieve]
\label{thm:sieve}
Let Assumptions~\ref{ass:sample} to \ref{ass:besov} hold with $p$, $q$ and the
constants fixed, and take
\begin{equation}
  \label{eq:Jn}
  J_n = \min\bigl\{J \in \mathbb{N} :\;
  2^{J} \ge (n/\log n)^{1/(2\seff+1)}\bigr\},
\end{equation}
so that $J_n = \log_2 n/(2\seff+1) + O(\log\log n)$. Then the requirement of
Proposition~\ref{prop:gram}(ii) holds along $J_n$, being itself of the exact
order $(\log n/n)^{2\seff/(2\seff+1)}$, and, with $\penn$ as in
\eqref{eq:lambda},
\[
  \normn{\widehat{f}-f}^{2}
  \;=\; O_p\Bigl\{\Bigl(\frac{\log n}{n}\Bigr)^{2\seff/(2\seff+1)}\Bigr\} .
\]
At this resolution the estimation term and the bias term are of the same
order, and the term $\sigma^{2}p/n$ of Theorem~\ref{thm:oracle}(iii) is of
strictly smaller order.
\end{theorem}

% derivations/05-taxas.tex, Corolario 4 (numeracao global).
\begin{corollary}[Components and coefficients]
\label{cor:components}
In the setting of Theorem~\ref{thm:sieve}, and with $r_n$ denoting the rate of
that theorem,
\[
  \sum_{\cv,\md}\bigl\lVert\widehat{g}_{\cv\md}-\gcomp{\cv}{\md}
  \bigr\rVert_{L_2[0,1]}^{2} = O_p(r_n),
  \qquad
  \bigl\lVert\widehat{\fcoef{\cv}}-\fcoef{\cv}
  \bigr\rVert_{L_2([0,1]^{q})} = O_p(r_n^{1/2}),
  \quad \cv = 1,\dots,p,
\]
and the same statements hold with the rate of Theorem~\ref{thm:main} in place
of $r_n$ under the conditions of that theorem.
\end{corollary}

The passage from the prediction norm to the norm of a single component is
where the design condition is used a second time, and the smallest eigenvalue
appears squared in the constant. Two further statements complete the picture.
The first dispenses with the design condition altogether, at the price of a
slower rate, and is the unconditional statement of this paper.

% derivations/05-taxas.tex, Proposicao 4 (numeracao global).
\begin{proposition}[Slow rate, no design condition]
\label{prop:slow}
Let Assumptions~\ref{ass:sample}, \ref{ass:basis} and \ref{ass:besov} hold,
with $\Amat$ of rank $p$ and $\penn$ as in \eqref{eq:lambda}. Then
$\normn{\widehat{f}-f}^{2} = O_p\{\penn\norm{\thetastar}_{1} +
\mathcal{B}_{J_n}^{2} + p/n\}$, where
$\norm{\thetastar}_{1} \le pqC_g\sum_{\lev<J_n}2^{\lev(1/2-s)}$ does not
depend on $\pi$. In particular, if $s > 1/2$ then
$\normn{\widehat{f}-f}^{2} = O_p\{(J_n/n)^{1/2}\}$ for every
$J_n \ge \log_2 n/(4\seff)$; and if $s < 1/2$, the choice
$2^{J_n} \asymp (n/J_n)^{1/(1-2s+4\seff)}$ gives
$\normn{\widehat{f}-f}^{2} = O_p\{(J_n/n)^{2\seff/(1-2s+4\seff)}\}$, an
exponent that equals $2s/(2s+1)$ when $\pi \ge 2$.
\end{proposition}

The second is the approximation lemma that produces Theorem~\ref{thm:main}.
For a vector $\btheta$ let $\theta_{(1)} \ge \theta_{(2)} \ge \cdots$ be the
nonincreasing rearrangement of the absolute values of its entries and let
$\btheta^{(M)}$ retain the $M$ largest of them.

% derivations/05-taxas.tex, Lema 9 (numeracao global).
\begin{lemma}[Besov regularity implies compressibility]
\label{lem:weak}
Under Assumptions~\ref{ass:basis} and \ref{ass:besov} with
$\pi(s+1/2) > 1$, the sieve oracle satisfies, for every $J$ and every
$k \ge 1$,
\[
  \theta^{*}_{(k)} \;\le\; C_\tau\,k^{-1/\tau},
  \qquad \tau = \frac{1}{s+1/2},
  \qquad C_\tau = C_g(pq\,A_{s,\pi})^{s+1/2},
  \qquad A_{s,\pi} = 1 + \frac{1}{1-2^{1-\pi(s+1/2)}},
\]
with constants that do not depend on $J$ and with $\tau \in (0,2)$ for every
$s > 0$. Consequently
$\sum_{k>M}\theta^{*2}_{(k)} \le \tau(2-\tau)^{-1}C_\tau^{2}M^{1-2/\tau}$ for
every $M \ge 1$.
\end{lemma}

Theorem~\ref{thm:main} follows by reading Theorem~\ref{thm:oracle}(ii) not at
the oracle but at the best $M$ term approximation of it, which is legitimate
because the proof of that theorem uses the comparator only through the
decomposition of $\bY$ and the size of its support. Balancing the two terms
that result, $\pen^{2}M/\gtil$ from estimation and
$M^{1-2/\tau}$ from the truncation, gives the effective dimension
\begin{equation}
  \label{eq:Mstar}
  \Mstar_n \;\asymp\; \Bigl(\frac{n}{J_n}\Bigr)^{\tau/2}
  \;\ll\; d_n \asymp pq\,2^{J_n},
\end{equation}
and a bound of order $\penn^{2-\tau}$, which is the rate of
Theorem~\ref{thm:main}. The three inequalities that define the window
\eqref{eq:window} are, in order, that the minimum in $M$ be interior, that the
bias of the sieve not dominate, and that the requirement of
Proposition~\ref{prop:gram}(ii) hold. Nothing in this argument says that the
estimator selects $\Mstar_n$ coordinates: no minimal signal or
irrepresentability condition is used, and what compressibility buys is a rate
governed by $\Mstar_n$ rather than by $d_n$, not the recovery of a support.

%%%%%%%%%%%%%%%%%%%%%%%%%%%%%%%%%%%%%%%%%%%%%%%%%%%%%%%%%%%%%%%%%%%%%%%%%%%%%%
\section{Computation and Tuning}
\label{sec:computation}

\subsection{The design matrix and the solver}
\label{sec:solver}

The design \eqref{eq:design} is built one block at a time: the $n \times \NJ$
matrix of the wavelets of $U_{\md}$ is evaluated once per modulator, by the
cascade algorithm of \citet{Daubechies-Lagarias-1991}, and then multiplied
rowwise by $X_{\cv}$. Since a compactly supported wavelet of level $\lev$ is
nonzero on an interval of length $(L-1)2^{-\lev}$, at most $L-1$ of the
$2^{\lev}$ columns of a level are nonzero at a given observation, so a block
has at most $n(L-1)J$ nonzero entries out of $n\NJ$ and the design is stored in
sparse form as soon as $J$ is moderate. Evaluating the basis once on a grid and
looking the values up removes the cascade from the inner loop of a
simulation.

Problem \eqref{eq:estimator} is then a lasso with $p$ unpenalized coordinates,
which coordinate descent solves along a path of penalty levels; we use the
\texttt{glmnet} package. Two points of scale deserve care, because they are
what ties the software to the theory. First, the penalty level reported for a
fit is the $\pen$ of \eqref{eq:estimator} and not the one of the solver, which
rescales its own by the ratio between the total number of columns and the
number of penalized ones when some coordinates carry zero penalty; the
conversion is checked by verifying the Karush-Kuhn-Tucker conditions of
\eqref{eq:estimator} at every point of the path. Second, when one linear
covariate is the constant, its column has zero variance, and the solver may
either drop it or share the level with its own intercept; the fit is the same
either way, but the level $\lvl{\cv}$ of the constant covariate is only
recovered after the intercept is added back to it.

\subsection{Rescaling and the boundary}
\label{sec:boundary}

The theory places $\bU$ in $[0,1]^q$ by assumption. In practice each modulator
is mapped to that interval by a monotone transformation, and, as in the
wavelet lasso of \citet*{Montoril-Pinheiro-Vidakovic-2019}, the image is
$[\epsilon, 1-\epsilon]$ with $\epsilon$ of the order $2^{-J}$, which keeps the
observations away from the strip where the periodization of
\eqref{eq:periodized} leaves an artifact; the price is that the fitted
component integrates to zero over the rescaled range and not over $[0,1]$, a
shift of level and not of shape. The alternative is the basis adapted to the
interval of \citet*{cohen1993wavelets}, which reproduces polynomials on $[0,1]$
and characterizes Besov regularity without any periodicity
\citep[Section 3.9]{Cohen-2003}, so that Lemma~\ref{lem:bias} holds with the
Besov space of the interval and the same proof. Its cost is structural: the
coarsest level must be large enough for the two boundary zones to be disjoint,
the constant is no longer a single column but a combination of the scaling
functions, and the constraint of \eqref{eq:additive} has to be imposed by
reparametrization, which leaves $pq(2^{\jzero}-1)$ unpenalized parameters
instead of none. We therefore keep the periodized basis as the default and
offer the interval basis as an option.

\subsection{Choosing the resolution and the penalty}
\label{sec:tuning}

The pair $(J,\pen)$ is chosen jointly. The default is $K$ fold cross
validation on a grid of resolutions, with the folds drawn once and reused for
every $J$, one design and one penalty path per resolution, and the usual two
readings of the curve: the minimizer of the cross validated error and the
largest penalty within one standard error of it. Two cheaper rules are also
available, the Bayesian information criterion and its extended version, with
degrees of freedom equal to the number of nonzero coefficients plus the $p$
levels, which is the usual unbiased count for the lasso.

The theory supplies a fourth rule, since \eqref{eq:Jn} and \eqref{eq:lambda}
are explicit. It is not a practical competitor as it stands, because it needs
$\sigma$, which has to be estimated, and $\seff$, which is unknown, and
because it is a statement about order, so that at moderate sample sizes it can
sit above the resolution that minimizes the realized error, the pre-asymptotic
regime being wide when $\seff$ is small. Its interest is as a reference for
what the cross validated choice ought to approach, and the difference between
the two is a finite sample quantity worth measuring. The four rules are
compared in the simulation study.

A remark on the structural variant. Replacing $\norm{\btheta}_1$ in
\eqref{eq:estimator} by the sparse group lasso penalty of
\citet{Simon-Friedman-Hastie-Tibshirani-2013}, with one group per block
$(\cv,\md)$, adds a second tuning parameter and removes entire components
rather than isolated coefficients, which is what a reader interested in
knowing which modulators act on which coefficients would want. It is a
different estimator, with a different oracle inequality
\citep{Lounici-Pontil-vandeGeer-Tsybakov-2011}, and the theory above does not
cover it; the choice between the two is settled empirically.

%%%%%%%%%%%%%%%%%%%%%%%%%%%%%%%%%%%%%%%%%%%%%%%%%%%%%%%%%%%%%%%%%%%%%%%%%%%%%%
\section*{Supplementary Materials}

The online Supplementary Material contains the proofs of all the results
stated in the paper, together with the auxiliary lemmas they use: the
identification argument and the properties of the periodized basis, the
approximation lemmas in $L_2$ and in the supremum norm and the proposition on
the cost of periodization, the population and empirical bounds for the Gram
matrix of the product design, the four lemmas behind the oracle inequality,
and the two lemmas behind the rates.

%%%%%%%%%%%%%%%%%%%%%%%%%%%%%%%%%%%%%%%%%%%%%%%%%%%%%%%%%%%%%%%%%%%%%%%%%%%%%%
\section*{Acknowledgements}

Write the acknowledgements here.
\par

%%%%%%%%%%%%%%%%%%%%%%%%%%%%%%%%%%%%%%%%%%%%%%%%%%%%%%%%%%%%%%%%%%%%%%%%%%%%%%

\bibhang=1.7pc
\bibsep=2pt
\fontsize{9}{14pt plus.8pt minus .6pt}\selectfont
\renewcommand\bibname{\large \bf References}
\expandafter\ifx\csname
natexlab\endcsname\relax\def\natexlab#1{#1}\fi
\expandafter\ifx\csname url\endcsname\relax
  \def\url#1{\texttt{#1}}\fi
\expandafter\ifx\csname urlprefix\endcsname\relax\def\urlprefix{URL}\fi

\bibliographystyle{chicago}
\bibliography{references_1}

%-------------------------------------------
\vskip .65cm
\noindent
first author affiliation
\vskip 2pt
\noindent
E-mail: (first author email)
\vskip 2pt

\noindent
second author affiliation
\vskip 2pt
\noindent
E-mail: (second author email)

\end{document}
